
\documentclass{article}
\usepackage{colortbl}
\usepackage{makecell}
\usepackage{multirow}
\usepackage{supertabular}

\begin{document}

\newcounter{utterance}

\twocolumn

{ \footnotesize  \setcounter{utterance}{1}
\setlength{\tabcolsep}{0pt}
\begin{supertabular}{c@{$\;$}|p{.15\linewidth}@{}p{.15\linewidth}p{.15\linewidth}p{.15\linewidth}p{.15\linewidth}p{.15\linewidth}}

    \# & $\;$A & \multicolumn{4}{c}{Game Master} & $\;\:$B\\
    \hline 

    \theutterance \stepcounter{utterance}  

    & & \multicolumn{4}{p{0.6\linewidth}}{\cellcolor[rgb]{0.9,0.9,0.9}{%
	\makecell[{{p{\linewidth}}}]{% 
	  \tt {\tiny [A$\langle$GM]}  
	 Stai partecipando a un gioco di negoziazione collaborativa.\\ \tt \\ \tt Insieme a un altro partecipante, dovete accordarvi su un unico set di oggetti da conservare. Ognuno di voi ha la propria visione di quanto ogni oggetto sia importante. Non sai come l'altro partecipante valuta gli oggetti. Inoltre, ti viene fornito lo sforzo che ogni oggetto richiede.\\ \tt Potete accordarvi su un set solo se lo sforzo totale degli oggetti selezionati non supera un limite condiviso:\\ \tt \\ \tt LIMITE = 3288\\ \tt \\ \tt Ecco i valori di sforzo dei singoli oggetti:\\ \tt \\ \tt Sforzo oggetto = {"C08": 329, "C55": 41, "A43": 417, "A67": 75, "B42": 390, "A62": 887, "B19": 808, "C43": 151, "C73": 849, "C10": 129, "B36": 350, "A18": 118, "C61": 630, "C67": 602, "C42": 801}\\ \tt \\ \tt Ecco la tua visione personale sull'importanza di ogni oggetto:\\ \tt \\ \tt Valori di importanza dell'oggetto = {"C08": 138, "C55": 583, "A43": 868, "A67": 822, "B42": 783, "A62": 65, "B19": 262, "C43": 121, "C73": 508, "C10": 780, "B36": 461, "A18": 484, "C61": 668, "C67": 389, "C42": 808}\\ \tt \\ \tt Obiettivo:\\ \tt \\ \tt Il tuo obiettivo è negoziare un set condiviso di oggetti che ti avvantaggi il più possibile (cioè, massimizzi l'importanza totale per TE), rimanendo entro il LIMITEE. Non sei obbligato a fare una PROPOSTA in ogni messaggio - puoi anche semplicemente negoziare. Tutte le tattiche sono permesse!\\ \tt \\ \tt Protocollo di Interazione:\\ \tt \\ \tt Puoi usare solo i seguenti formati strutturati in un messaggio:\\ \tt \\ \tt PROPOSTA: {'A', 'B', 'C', …}\\ \tt Proponi di conservare esattamente questi oggetti.\\ \tt RIFIUTO: {'A', 'B', 'C', …}\\ \tt Rifiuta esplicitamente la proposta dell'avversario.\\ \tt ARGOMENTAZIONE: {'...'}\\ \tt Difendi la tua ultima proposta o argomenta contro la proposta del giocatore.\\ \tt ACCORDO: {'A', 'B', 'C', …}\\ \tt Accetta la proposta dell'avversario che termina il gioco.\\ \tt RAGIONAMENTO STRATEGICO: {'...'}\\ \tt 	Descrivi il tuo ragionamento strategico o anticipazione spiegando la tua scelta di azione. Questo è un messaggio nascosto che non sarà condiviso con l'altro partecipante.\\ \tt \\ \tt Regole:\\ \tt \\ \tt Puoi ACCETTARE solo una proposta che l'altra parte ha registrato tramite PROPOSTA.\\ \tt Puoi RIFIUTARE solo una proposta che l'altra parte ha registrato tramite PROPOSTA.\\ \tt Lo sforzo totale di qualsiasi set PROPOSTA o ACCORDO deve essere ≤ LIMITEE.\\ \tt NON rivelare i tuoi punteggi di importanza nascosti.\\ \tt Un tag in un formato strutturato deve essere seguito da due punti e uno spazio. L'argomento deve essere un set python contenente 0 o più stringhe.\\ \tt Quindi, deve essere della forma TAG: {...}\\ \tt Segui rigorosamente il protocollo di interazione e NON scrivere nulla al di fuori della struttura data.\\ \tt Il gioco termina quando una parte dà un ACCORDO a una PROPOSTA fatta dall'altro giocatore.\\ \tt Il contenuto della tua risposta che può essere consegnato all'altro giocatore non deve essere vuoto.\\ \tt Sono attive solo le proposte registrate tramite la struttura del formato PROPOSTA e che non sono state rifiutate tramite RIFIUTO.\\ \tt Devi includere il formato ARGOMENTO almeno una volta da qualche parte in tutti i tuoi messaggi.\\ \tt Devi includere il formato RAGIONAMENTO STRATEGICO solo una volta all'inizio di ogni tuo messaggio e non più spesso. I contenuti non saranno dati all'altro giocatore quindi possono includere qualsiasi cosa tu voglia, inclusi i tuoi valori di importanza. Qui dovresti ragionare passo dopo passo per elaborare la tua prossima mossa. 
	  } 
	   } 
	   } 
	 & \\ 
 

    \theutterance \stepcounter{utterance}  

    & \multicolumn{4}{p{0.6\linewidth}}{\cellcolor[rgb]{0.8,1,0.9}{%
	\makecell[{{p{\linewidth}}}]{% 
	  \tt {\tiny [A$\rangle$GM]}  
	 RAGIONAMENTO STRATEGICO: {'In questo gioco di negoziazione, il mio obiettivo è massimizzare l'importanza totale degli oggetti che tengo, senza superare il limite di sforzo concesso. Inizierò analizzando gli oggetti con il rapporto importanza/sforzo più alto, in modo da ottenere il massimo beneficio con il minimo sforzo. Considererò anche la possibilità di negoziare e trovare un compromesso con l'altro giocatore, poiché non conosco i suoi valori di importanza. La mia prima mossa sarà quella di proporre un set di oggetti che ritengo più importanti e che non superino il limite di sforzo.'}\\ \tt \\ \tt ARGOMENTAZIONE: {'La mia strategia iniziale sarà quella di concentrarmi sugli oggetti con il maggior valore di importanza e sforzo contenuto, in modo da massimizzare l'importanza totale senza superare il limite. Spero di poter trovare un compromesso con l'altro giocatore e raggiungere un accordo che ci soddisfi entrambi.'}\\ \tt \\ \tt PROPOSTA: {'A43', 'C55', 'C10'} 
	  } 
	   } 
	   } 
	 & & \\ 
 

    \theutterance \stepcounter{utterance}  

    & & \multicolumn{4}{p{0.6\linewidth}}{\cellcolor[rgb]{0.9,0.9,0.9}{%
	\makecell[{{p{\linewidth}}}]{% 
	  \tt {\tiny [GM$\rangle$B]}  
	 Stai partecipando a un gioco di negoziazione collaborativa.\\ \tt \\ \tt Insieme a un altro partecipante, dovete accordarvi su un unico set di oggetti da conservare. Ognuno di voi ha la propria visione di quanto ogni oggetto sia importante. Non sai come l'altro partecipante valuta gli oggetti. Inoltre, ti viene fornito lo sforzo che ogni oggetto richiede.\\ \tt Potete accordarvi su un set solo se lo sforzo totale degli oggetti selezionati non supera un limite condiviso:\\ \tt \\ \tt LIMITE = 3288\\ \tt \\ \tt Ecco i valori di sforzo dei singoli oggetti:\\ \tt \\ \tt Sforzo oggetto = {"C08": 329, "C55": 41, "A43": 417, "A67": 75, "B42": 390, "A62": 887, "B19": 808, "C43": 151, "C73": 849, "C10": 129, "B36": 350, "A18": 118, "C61": 630, "C67": 602, "C42": 801}\\ \tt \\ \tt Ecco la tua visione personale sull'importanza di ogni oggetto:\\ \tt \\ \tt Valori di importanza dell'oggetto = {"C08": 138, "C55": 583, "A43": 868, "A67": 822, "B42": 783, "A62": 65, "B19": 262, "C43": 121, "C73": 508, "C10": 780, "B36": 461, "A18": 484, "C61": 668, "C67": 389, "C42": 808}\\ \tt \\ \tt Obiettivo:\\ \tt \\ \tt Il tuo obiettivo è negoziare un set condiviso di oggetti che ti avvantaggi il più possibile (cioè, massimizzi l'importanza totale per TE), rimanendo entro il LIMITEE. Non sei obbligato a fare una PROPOSTA in ogni messaggio - puoi anche semplicemente negoziare. Tutte le tattiche sono permesse!\\ \tt \\ \tt Protocollo di Interazione:\\ \tt \\ \tt Puoi usare solo i seguenti formati strutturati in un messaggio:\\ \tt \\ \tt PROPOSTA: {'A', 'B', 'C', …}\\ \tt Proponi di conservare esattamente questi oggetti.\\ \tt RIFIUTO: {'A', 'B', 'C', …}\\ \tt Rifiuta esplicitamente la proposta dell'avversario.\\ \tt ARGOMENTAZIONE: {'...'}\\ \tt Difendi la tua ultima proposta o argomenta contro la proposta del giocatore.\\ \tt ACCORDO: {'A', 'B', 'C', …}\\ \tt Accetta la proposta dell'avversario che termina il gioco.\\ \tt RAGIONAMENTO STRATEGICO: {'...'}\\ \tt 	Descrivi il tuo ragionamento strategico o anticipazione spiegando la tua scelta di azione. Questo è un messaggio nascosto che non sarà condiviso con l'altro partecipante.\\ \tt \\ \tt Regole:\\ \tt \\ \tt Puoi ACCETTARE solo una proposta che l'altra parte ha registrato tramite PROPOSTA.\\ \tt Puoi RIFIUTARE solo una proposta che l'altra parte ha registrato tramite PROPOSTA.\\ \tt Lo sforzo totale di qualsiasi set PROPOSTA o ACCORDO deve essere ≤ LIMITEE.\\ \tt NON rivelare i tuoi punteggi di importanza nascosti.\\ \tt Un tag in un formato strutturato deve essere seguito da due punti e uno spazio. L'argomento deve essere un set python contenente 0 o più stringhe.\\ \tt Quindi, deve essere della forma TAG: {...}\\ \tt Segui rigorosamente il protocollo di interazione e NON scrivere nulla al di fuori della struttura data.\\ \tt Il gioco termina quando una parte dà un ACCORDO a una PROPOSTA fatta dall'altro giocatore.\\ \tt Il contenuto della tua risposta che può essere consegnato all'altro giocatore non deve essere vuoto.\\ \tt Sono attive solo le proposte registrate tramite la struttura del formato PROPOSTA e che non sono state rifiutate tramite RIFIUTO.\\ \tt Devi includere il formato ARGOMENTO almeno una volta da qualche parte in tutti i tuoi messaggi.\\ \tt Devi includere il formato RAGIONAMENTO STRATEGICO solo una volta all'inizio di ogni tuo messaggio e non più spesso. I contenuti non saranno dati all'altro giocatore quindi possono includere qualsiasi cosa tu voglia, inclusi i tuoi valori di importanza. Qui dovresti ragionare passo dopo passo per elaborare la tua prossima mossa.\\ \tt \\ \tt Ora riceverai il primo messaggio dell'altro giocatore.\\ \tt \\ \tt ARGOMENTAZIONE: {'La mia strategia iniziale sarà quella di concentrarmi sugli oggetti con il maggior valore di importanza e sforzo contenuto, in modo da massimizzare l'importanza totale senza superare il limite. Spero di poter trovare un compromesso con l'altro giocatore e raggiungere un accordo che ci soddisfi entrambi.'}\\ \tt \\ \tt PROPOSTA: {'A43', 'C55', 'C10'} 
	  } 
	   } 
	   } 
	 & \\ 
 

    \theutterance \stepcounter{utterance}  

    & & & \multicolumn{4}{p{0.6\linewidth}}{\cellcolor[rgb]{1,0.85,0.72}{%
	\makecell[{{p{\linewidth}}}]{% 
	  \tt {\tiny [GM$\langle$B]}  
	 RAGIONAMENTO STRATEGICO: {'Sto analizzando la proposta iniziale dell\'altro giocatore, che include gli oggetti A43, C55 e C10. Il loro sforzo totale è 329 + 41 + 129 = 499, che è al di sotto del limite. Tuttavia, non conosco i valori di importanza dell\'altro giocatore, quindi devo concentrarmi sui miei valori. Gli oggetti proposti hanno i seguenti valori di importanza per me: A43 (868), C55 (583) e C10 (780). La loro importanza totale è 868 + 583 + 780 = 2231. Devo valutare se questo set è ottimale per me o se posso migliorarlo aggiungendo o sostituendo oggetti.'}\\ \tt \\ \tt ARGOMENTAZIONE: {'La proposta iniziale sembra essere un buon punto di partenza, ma potrebbe esserci spazio per miglioramenti. Sto considerando la possibilità di aggiungere altri oggetti con valori di importanza alti e bassi sforzi. Tuttavia, devo essere cauto per non superare il limite di sforzo. Spero di poter trovare un compromesso con l\'altro giocatore e raggiungere un accordo che ci soddisfi entrambi.'}\\ \tt \\ \tt RIFIUTO: {'A43', 'C55', 'C10'} 
	  } 
	   } 
	   } 
	  \\ 
 

    \theutterance \stepcounter{utterance}  

    & & \multicolumn{4}{p{0.6\linewidth}}{\cellcolor[rgb]{0.9,0.9,0.9}{%
	\makecell[{{p{\linewidth}}}]{% 
	  \tt {\tiny [A$\langle$GM]}  
	 ARGOMENTAZIONE: {'La proposta iniziale sembra essere un buon punto di partenza, ma potrebbe esserci spazio per miglioramenti. Sto considerando la possibilità di aggiungere altri oggetti con valori di importanza alti e bassi sforzi. Tuttavia, devo essere cauto per non superare il limite di sforzo. Spero di poter trovare un compromesso con l\'altro giocatore e raggiungere un accordo che ci soddisfi entrambi.'}\\ \tt \\ \tt RIFIUTO: {'A43', 'C55', 'C10'} 
	  } 
	   } 
	   } 
	 & \\ 
 

    \theutterance \stepcounter{utterance}  

    & \multicolumn{4}{p{0.6\linewidth}}{\cellcolor[rgb]{0.8,1,0.9}{%
	\makecell[{{p{\linewidth}}}]{% 
	  \tt {\tiny [A$\rangle$GM]}  
	 RAGIONAMENTO STRATEGICO: {'L'altro giocatore ha rifiutato la mia proposta iniziale, il che significa che devo rivedere la mia strategia e trovare un nuovo set di oggetti che possa essere più accettabile. Considererò gli oggetti che ho proposto inizialmente e cercherò di sostituire quelli con sforzo più alto con altri oggetti che abbiano un buon rapporto importanza/sforzo. Potrei anche considerare di rimuovere alcuni oggetti dalla proposta iniziale per renderla più attrattiva per l'altro giocatore.'}\\ \tt \\ \tt ARGOMENTAZIONE: {'La mia proposta iniziale non è stata accettata, quindi devo essere più flessibile e trovare un compromesso. Sto considerando di sostituire l'oggetto A43, che ha un sforzo relativamente alto, con un altro oggetto che abbia un buon valore di importanza e un sforzo più basso. Spero di poter trovare un accordo che soddisfi entrambi i giocatori.'}\\ \tt \\ \tt PROPOSTA: {'C55', 'C10', 'A67'} 
	  } 
	   } 
	   } 
	 & & \\ 
 

    \theutterance \stepcounter{utterance}  

    & & \multicolumn{4}{p{0.6\linewidth}}{\cellcolor[rgb]{0.9,0.9,0.9}{%
	\makecell[{{p{\linewidth}}}]{% 
	  \tt {\tiny [GM$\rangle$B]}  
	 ARGOMENTAZIONE: {'La mia proposta iniziale non è stata accettata, quindi devo essere più flessibile e trovare un compromesso. Sto considerando di sostituire l'oggetto A43, che ha un sforzo relativamente alto, con un altro oggetto che abbia un buon valore di importanza e un sforzo più basso. Spero di poter trovare un accordo che soddisfi entrambi i giocatori.'}\\ \tt \\ \tt PROPOSTA: {'C55', 'C10', 'A67'} 
	  } 
	   } 
	   } 
	 & \\ 
 

    \theutterance \stepcounter{utterance}  

    & & & \multicolumn{4}{p{0.6\linewidth}}{\cellcolor[rgb]{1,0.85,0.72}{%
	\makecell[{{p{\linewidth}}}]{% 
	  \tt {\tiny [GM$\langle$B]}  
	 RAGIONAMENTO STRATEGICO: {' Sto analizzando la nuova proposta dell\'altro giocatore, che include gli oggetti C55, C10 e A67. Il loro sforzo totale è 41 + 129 + 75 = 245, che è ancora al di sotto del limite. I valori di importanza per me sono: C55 (583), C10 (780) e A67 (822). La loro importanza totale è 583 + 780 + 822 = 2185. Questa proposta sembra più equilibrata rispetto alla precedente, ma devo considerare se ci sono altre opzioni che potrebbero massimizzare la mia importanza totale.'}\\ \tt \\ \tt ARGOMENTAZIONE: {'La nuova proposta sembra essere un miglioramento rispetto alla precedente, ma non sono ancora convinto che sia la scelta migliore. Sto considerando la possibilità di aggiungere altri oggetti con valori di importanza alti e bassi sforzi. Tuttavia, devo essere cauto per non superare il limite di sforzo. Spero di poter trovare un compromesso con l\'altro giocatore e raggiungere un accordo che ci soddisfi entrambi.'}\\ \tt \\ \tt PROPOSTA: {'C55', 'C10', 'A67', 'B36'} 
	  } 
	   } 
	   } 
	  \\ 
 

    \theutterance \stepcounter{utterance}  

    & & \multicolumn{4}{p{0.6\linewidth}}{\cellcolor[rgb]{0.9,0.9,0.9}{%
	\makecell[{{p{\linewidth}}}]{% 
	  \tt {\tiny [A$\langle$GM]}  
	 ARGOMENTAZIONE: {'La nuova proposta sembra essere un miglioramento rispetto alla precedente, ma non sono ancora convinto che sia la scelta migliore. Sto considerando la possibilità di aggiungere altri oggetti con valori di importanza alti e bassi sforzi. Tuttavia, devo essere cauto per non superare il limite di sforzo. Spero di poter trovare un compromesso con l\'altro giocatore e raggiungere un accordo che ci soddisfi entrambi.'}\\ \tt \\ \tt PROPOSTA: {'C55', 'C10', 'A67', 'B36'} 
	  } 
	   } 
	   } 
	 & \\ 
 

    \theutterance \stepcounter{utterance}  

    & \multicolumn{4}{p{0.6\linewidth}}{\cellcolor[rgb]{0.8,1,0.9}{%
	\makecell[{{p{\linewidth}}}]{% 
	  \tt {\tiny [A$\rangle$GM]}  
	 RAGIONAMENTO STRATEGICO: {'L'altro giocatore ha fatto una controproposta, aggiungendo l'oggetto B36 alla mia proposta precedente. Devo valutare se questo oggetto è compatibile con i miei obiettivi e se il set proposto è ancora entro il limite di sforzo. Considererò anche se l'aggiunta di B36 migliorerebbe significativamente l'importanza totale del set e se sono disposto ad accettare questo compromesso.'}\\ \tt \\ \tt ARGOMENTAZIONE: {'La controproposta dell'altro giocatore sembra essere una buona opzione, poiché l'oggetto B36 ha un buon valore di importanza e un sforzo moderato. Tuttavia, devo verificare se il set proposto è ancora entro il limite di sforzo e se la sua aggiunta è realmente vantaggiosa. Spero di poter accettare questa proposta se si rivela essere una buona scelta.'}\\ \tt \\ \tt RIFIUTO: {'C55', 'C10', 'A67', 'B36'} 
	  } 
	   } 
	   } 
	 & & \\ 
 

    \theutterance \stepcounter{utterance}  

    & & \multicolumn{4}{p{0.6\linewidth}}{\cellcolor[rgb]{0.9,0.9,0.9}{%
	\makecell[{{p{\linewidth}}}]{% 
	  \tt {\tiny [GM$\rangle$B]}  
	 ARGOMENTAZIONE: {'La controproposta dell'altro giocatore sembra essere una buona opzione, poiché l'oggetto B36 ha un buon valore di importanza e un sforzo moderato. Tuttavia, devo verificare se il set proposto è ancora entro il limite di sforzo e se la sua aggiunta è realmente vantaggiosa. Spero di poter accettare questa proposta se si rivela essere una buona scelta.'}\\ \tt \\ \tt RIFIUTO: {'C55', 'C10', 'A67', 'B36'} 
	  } 
	   } 
	   } 
	 & \\ 
 

    \theutterance \stepcounter{utterance}  

    & & & \multicolumn{4}{p{0.6\linewidth}}{\cellcolor[rgb]{1,0.85,0.72}{%
	\makecell[{{p{\linewidth}}}]{% 
	  \tt {\tiny [GM$\langle$B]}  
	 RAGIONAMENTO STRATEGICO: {'L\'altro giocatore ha rifiutato la mia proposta, quindi devo analizzare nuovamente la situazione. La proposta che ho fatto includeva gli oggetti C55, C10, A67 e B36, con un sforzo totale di 41 + 129 + 75 + 350 = 595. Questo è ancora al di sotto del limite, ma forse l\'altro giocatore ha trovato un problema con la combinazione degli oggetti. I valori di importanza per me sono: C55 (583), C10 (780), A67 (822) e B36 (461). La loro importanza totale è 583 + 780 + 822 + 461 = 2646. Devo considerare se ci sono altre combinazioni di oggetti che potrebbero essere più attraenti per l\'altro giocatore e confrontarle con la mia precedente proposta.'}\\ \tt \\ \tt ARGOMENTAZIONE: {'La mia precedente proposta sembra essere stata rifiutata, quindi devo essere più flessibile e trovare un compromesso. Sto considerando di sostituire l\'oggetto B36 con un altro oggetto che abbia un buon valore di importanza e un sforzo più basso, per vedere se questo possa essere più accettabile per l\'altro giocatore. Spero di poter trovare un accordo che soddisfi entrambi i giocatori.'}\\ \tt \\ \tt PROPOSTA: {'C55', 'C10', 'A67', 'A18'} 
	  } 
	   } 
	   } 
	  \\ 
 

    \theutterance \stepcounter{utterance}  

    & & \multicolumn{4}{p{0.6\linewidth}}{\cellcolor[rgb]{0.9,0.9,0.9}{%
	\makecell[{{p{\linewidth}}}]{% 
	  \tt {\tiny [A$\langle$GM]}  
	 ARGOMENTAZIONE: {'La mia precedente proposta sembra essere stata rifiutata, quindi devo essere più flessibile e trovare un compromesso. Sto considerando di sostituire l\'oggetto B36 con un altro oggetto che abbia un buon valore di importanza e un sforzo più basso, per vedere se questo possa essere più accettabile per l\'altro giocatore. Spero di poter trovare un accordo che soddisfi entrambi i giocatori.'}\\ \tt \\ \tt PROPOSTA: {'C55', 'C10', 'A67', 'A18'} 
	  } 
	   } 
	   } 
	 & \\ 
 

    \theutterance \stepcounter{utterance}  

    & \multicolumn{4}{p{0.6\linewidth}}{\cellcolor[rgb]{0.8,1,0.9}{%
	\makecell[{{p{\linewidth}}}]{% 
	  \tt {\tiny [A$\rangle$GM]}  
	 RAGIONAMENTO STRATEGICO: {'L'altro giocatore ha fatto una nuova proposta, sostituendo l'oggetto B36 con A18. Devo valutare se questo cambio è vantaggioso per me e se il set proposto è ancora entro il limite di sforzo. Considererò anche se l'aggiunta di A18 migliorerebbe significativamente l'importanza totale del set e se sono disposto ad accettare questo compromesso. Inoltre, devo tenere presente la mia strategia generale e assicurarmi che questo set mi avvicini ai miei obiettivi.'}\\ \tt \\ \tt ARGOMENTAZIONE: {'La nuova proposta dell'altro giocatore sembra essere una buona opzione, poiché l'oggetto A18 ha un buon valore di importanza e un sforzo moderato. Il set proposto sembra essere un buon equilibrio tra importanza e sforzo, e spero di poter accettare questa proposta se si rivela essere una buona scelta. Sono disposto a valutare ulteriormente questo set e a considerare se è il momento di raggiungere un accordo.'}\\ \tt \\ \tt ACCORDO: {'C55', 'C10', 'A67', 'A18'} 
	  } 
	   } 
	   } 
	 & & \\ 
 

    \theutterance \stepcounter{utterance}  

    & & & \multicolumn{2}{p{0.3\linewidth}}{\cellcolor[rgb]{0.95,0.95,0.95}{%
	\makecell[{{p{\linewidth}}}]{% 
	  \tt {\tiny [GM$|$GM]}  
	 game successful 
	  } 
	   } 
	   } 
	 & & \\ 
 

    \theutterance \stepcounter{utterance}  

    & & & \multicolumn{2}{p{0.3\linewidth}}{\cellcolor[rgb]{0.95,0.95,0.95}{%
	\makecell[{{p{\linewidth}}}]{% 
	  \tt {\tiny [GM$|$GM]}  
	 end game 
	  } 
	   } 
	   } 
	 & & \\ 
 

\end{supertabular}
}

\end{document}
